\chapter{Why Single Prompts Fail}

LLMs have fundamental constraints:

\begin{itemize}
    \item \textbf{Context windows}: Limited memory for each conversation
    \item \textbf{Statelessness}: No persistent understanding across sessions
    \item \textbf{Hallucination tendency}: Confident generation without knowledge
    \item \textbf{Global awareness limitations}: Poor integration of distant parts of long inputs
\end{itemize}

These constraints mean that complex tasks cannot be solved with single prompts. The input is too messy, the reasoning too long, the verification too incomplete.

\textbf{Prompt chaining} addresses these constraints by decomposing complex cognition into manageable steps:

\begin{lstlisting}
Step 1: Extract structure from raw input
Step 2: Identify gaps and ambiguities
Step 3: Clarify each component separately
Step 4: Synthesize into coherent whole
Step 5: Verify against requirements
Step 6: Iterate on failures
\end{lstlisting}

Each step is a discrete prompt with a focused objective. The chain creates a \textbf{cognitive pipeline}---a series of transformations that convert chaos into clarity.

\begin{quotebox}
\textbf{Agentic Era Reality:} Claude Code performs this chaining automatically:
\begin{itemize}
    \item The \textbf{Task tool} spawns specialized agents for different chain steps
    \item \textbf{Plan Mode} implements research→plan→implement chains
    \item \textbf{Extended thinking} handles internal reasoning chains
    \item \textbf{Tool orchestration} chains file reads, edits, and verification
\end{itemize}
Your role shifts from \emph{building chains} to \emph{evaluating chain outputs} and \emph{intervening when chains fail}.
\end{quotebox}
